\section{Problema}

\begin{frame}{Duas saídas não significam apenas dois testes}
  \begin{columns}[T,onlytextwidth]
    \begin{column}{.38\textwidth}
      \begin{block}{Saída}
        \centering
        \textbf{true}\\
        acesso concedido
      \end{block}
      \begin{block}{Saída}
        \centering
        \textbf{false}\\
        acesso negado
      \end{block}
    \end{column}
    \begin{column}{.56\textwidth}
      Usuário e senha pertencem a classes diferentes:
      \begin{itemize}
        \item usuário ativo, bloqueado, inexistente ou vazio;
        \item senha correspondente, incorreta, alterada ou vazia;
        \item cada combinação pode exigir mensagem e estado diferentes.
      \end{itemize}
    \end{column}
  \end{columns}
\end{frame}

\begin{frame}{Objetivos da atividade}
  Ao final, sua equipe deverá ser capaz de:
  \begin{itemize}
    \item transformar requisitos em classes de equivalência;
    \item distinguir classe de entrada de resultado do teste;
    \item combinar classes sem testar todas as possibilidades;
    \item elaborar plano, matriz de partições e casos de teste;
    \item executar os casos no LoginLab;
    \item registrar observado, evidência, defeito e decisão.
  \end{itemize}
\end{frame}

\begin{frame}{Sistema sob teste: LoginLab 1.0.0}
  Abra \texttt{index.html}. A massa conhecida contém:
  \vspace{4mm}

  \small
  \begin{tabular}{p{36mm}p{37mm}p{27mm}}
    \toprule
    \textbf{Usuário} & \textbf{Senha} & \textbf{Estado}\\
    \midrule
    \texttt{aluno01} & \texttt{Teste@123} & ativo\\
    \texttt{professor01} & \texttt{Aula@2026} & ativo\\
    \texttt{bloqueado01} & \texttt{Segura@456} & bloqueado\\
    \bottomrule
  \end{tabular}
  \vspace{6mm}

  \begin{alertblock}{Importante}
    Use somente as contas didáticas. O microprojeto é local e não utiliza dados reais.
  \end{alertblock}
\end{frame}

\section{Requisitos}

\begin{frame}{Regras de autenticação}
  \small
  \begin{tabular}{p{24mm}p{91mm}}
    \toprule
    \textbf{ID} & \textbf{Regra}\\
    \midrule
    RN-LOGIN-01 & Usuário e senha são obrigatórios.\\
    RN-LOGIN-02 & Usuário ativo com senha correspondente deve autenticar.\\
    RN-LOGIN-03 & Senha incorreta não deve autenticar.\\
    RN-LOGIN-04 & Usuário inexistente não deve autenticar.\\
    RN-LOGIN-05 & Usuário bloqueado não autentica, mesmo com senha correta.\\
    \bottomrule
  \end{tabular}
\end{frame}

\begin{frame}{Mensagem e sessão também fazem parte do resultado}
  \small
  \begin{tabular}{p{24mm}p{91mm}}
    \toprule
    \textbf{ID} & \textbf{Regra}\\
    \midrule
    RN-LOGIN-06 & Credencial inválida exibe
      \texttt{Usuário ou senha inválidos}.\\
    RN-LOGIN-07 & Conta bloqueada exibe
      \texttt{Conta bloqueada. Procure o suporte}.\\
    RN-LOGIN-08 & A senha diferencia maiúsculas e minúsculas.\\
    RN-LOGIN-09 & Falha de autenticação não cria sessão.\\
    \bottomrule
  \end{tabular}
  \vspace{5mm}

  \begin{block}{Oráculo}
    Não observe apenas true ou false: compare mensagem e estado da sessão.
  \end{block}
\end{frame}

\section{Artefatos}

\begin{frame}{Parte 1 — produzir os artefatos antes de testar}
  \begin{enumerate}
    \item elaborar o plano de teste;
    \item identificar classes de equivalência;
    \item selecionar combinações representativas;
    \item escrever de 8 a 10 casos de teste.
  \end{enumerate}
  \vspace{5mm}

  \begin{alertblock}{Disciplina de execução}
    Resultado real e status permanecem vazios até o caso ser executado.
  \end{alertblock}
\end{frame}

\begin{frame}{Artefato 1 — plano de teste}
  Use \texttt{modelos/plano\_de\_teste.md} para definir:
  \begin{itemize}
    \item objetivo e decisão apoiada;
    \item escopo e fora de escopo;
    \item riscos e prioridades;
    \item estratégia e técnica;
    \item ambiente e massa de dados;
    \item critérios de entrada, saída, suspensão e retomada.
  \end{itemize}
\end{frame}

\begin{frame}{Artefato 2 — classes do campo usuário}
  \small
  \begin{tabular}{p{34mm}p{39mm}p{37mm}}
    \toprule
    \textbf{Classe} & \textbf{Representante} & \textbf{Tratamento}\\
    \midrule
    cadastrado e ativo & \texttt{aluno01} & pode autenticar\\
    cadastrado e bloqueado & \texttt{bloqueado01} & negar\\
    inexistente & \texttt{fantasma01} & negar\\
    vazio & campo sem valor & validar obrigatoriedade\\
    \bottomrule
  \end{tabular}
  \vspace{5mm}

  \begin{block}{Pergunta}
    Usuário cadastrado é sempre uma entrada válida para autenticação?
  \end{block}
\end{frame}

\begin{frame}{Artefato 2 — classes do campo senha}
  \small
  \begin{tabular}{p{41mm}p{37mm}p{32mm}}
    \toprule
    \textbf{Classe} & \textbf{Representante} & \textbf{Tratamento}\\
    \midrule
    correspondente ao usuário & \texttt{Teste@123} & pode autenticar\\
    incorreta com formato válido & \texttt{Errada@123} & negar\\
    diferença de caixa & \texttt{teste@123} & negar\\
    vazia & campo sem valor & validar obrigatoriedade\\
    \bottomrule
  \end{tabular}
  \vspace{5mm}

  \begin{alertblock}{Dependência}
    ``Senha correta'' depende do usuário informado. A mesma senha pode ser correta
    para uma conta e incorreta para outra.
  \end{alertblock}
\end{frame}

\begin{frame}{Artefato 3 — selecionar combinações}
  \small
  \begin{tabular}{p{38mm}p{39mm}p{24mm}}
    \toprule
    \textbf{Usuário} & \textbf{Senha} & \textbf{Resultado}\\
    \midrule
    ativo & correspondente & true\\
    ativo & incorreta & false\\
    inexistente & qualquer não vazia & false\\
    bloqueado & correspondente & false\\
    vazio & qualquer & false\\
    \bottomrule
  \end{tabular}
  \vspace{4mm}

  Selecione combinações que:
  \begin{itemize}
    \item cubram cada classe pelo menos uma vez;
    \item isolem o motivo da rejeição;
    \item priorizem segurança e evitem redundância.
  \end{itemize}
\end{frame}

\begin{frame}{Artefato 4 — especificar de 8 a 10 casos}
  Cada caso deve conter:
  \begin{itemize}
    \item ID e título;
    \item requisitos e riscos relacionados;
    \item prioridade e pré-condições;
    \item usuário e senha concretos;
    \item passos necessários;
    \item mensagem e estado de sessão esperados.
  \end{itemize}
  \begin{block}{Modelo}
    Use \texttt{modelos/casos\_de\_teste.md}.
  \end{block}
\end{frame}

\section{Execução}

\begin{frame}{Parte 2 — executar o microprojeto}
  Para cada caso:
  \begin{enumerate}
    \item selecione \textbf{Limpar sessão};
    \item informe usuário e senha;
    \item selecione \textbf{Entrar};
    \item registre mensagem, estado da sessão e status;
    \item capture evidência se houver divergência.
  \end{enumerate}
  \vspace{4mm}

  Use \texttt{modelos/registro\_de\_execucao.md}.
\end{frame}

\begin{frame}{Parte 3 — interpretar os resultados}
  \small
  \begin{enumerate}
    \item Quantas classes de usuário foram identificadas?
    \item Quantas classes de senha foram identificadas?
    \item Por que não bastam um caso true e um caso false?
    \item Quais combinações diferentes produziram false?
    \item A mensagem revela se um usuário está cadastrado?
    \item Alguma divergência deve bloquear a liberação?
  \end{enumerate}
  \vspace{3mm}

  Se houver falha, preencha \texttt{modelos/defeito.md}.
\end{frame}

\begin{frame}{Entregáveis}
  \begin{columns}[T,onlytextwidth]
    \begin{column}{.48\textwidth}
      \begin{enumerate}
        \item plano de teste;
        \item classes de equivalência;
        \item matriz de combinações;
        \item 8 a 10 casos.
      \end{enumerate}
    \end{column}
    \begin{column}{.48\textwidth}
      \begin{enumerate}
        \setcounter{enumi}{4}
        \item registro de execução;
        \item evidências;
        \item relatório de defeito;
        \item decisão de liberação.
      \end{enumerate}
    \end{column}
  \end{columns}
\end{frame}

\begin{frame}{Tempo sugerido}
  \small
  \begin{tabular}{p{35mm}p{72mm}}
    \toprule
    \textbf{Tempo} & \textbf{Atividade}\\
    \midrule
    20 minutos & requisitos, riscos e plano\\
    20 minutos & partições e combinações\\
    20 minutos & casos de teste\\
    25 minutos & execução e evidências\\
    15 minutos & revisão e decisão\\
    \bottomrule
  \end{tabular}
\end{frame}

\begin{frame}{Critérios de avaliação}
  \small
  \begin{tabular}{p{91mm}r}
    \toprule
    \textbf{Critério} & \textbf{Peso}\\
    \midrule
    Classes derivadas dos requisitos & 25\%\\
    Seleção não redundante de combinações & 20\%\\
    Clareza e reprodutibilidade dos casos & 20\%\\
    Registro do observado e evidências & 20\%\\
    Rastreabilidade e decisão final & 15\%\\
    \bottomrule
  \end{tabular}
\end{frame}

\begin{frame}{A pergunta que encerra a atividade}
  \centering
  \vfill
  {\LARGE\color{azulinstitucional}\bfseries
    Se o resultado é apenas true ou false,\\[3mm]
    por que precisamos de vários casos?}
  \vfill
  \begin{block}{Ponto de partida}
    Abra \texttt{index.html}. Preencha os artefatos antes de executar.
  \end{block}
\end{frame}

